% Template pour mémoires et thèses - Université Laval
% Adapté pour la conformité aux normes de la FESP (Faculté des études
% supérieures et postdoctorales), mise à jour novembre 2025.
%
% Basé sur le Reed College LaTeX thesis template par Sam Noble (SN),
% Ben Salzberg (BTS), et Jess Youngberg (JY).
%
% Adaptation pour l'Université Laval par Adrien Cloutier (2025):
% - Conformité aux normes FESP (marges 30mm/25mm, interligne 1.5)
% - Support pour format par articles
% - Avant-propos avec déclaration IA obligatoire (2026)
% - Copyright dynamique avec année courante
% - Élimination des pages blanches
%
% Original Reed College template: http://web.reed.edu/cis/help/latex.html
% Documentation FESP: https://www.fesp.ulaval.ca/memoires-et-theses/normes-de-presentation
%
% Toute ligne commençant par % est un commentaire et n'apparaît pas dans le document.
% Très utile pour le dépannage et les notes personnelles.

%%
%% Preamble
%%
% \documentclass{<something>} must begin each LaTeX document
\documentclass[12pt]{ulaval}
% Packages are extensions to the basic LaTeX functions. Whatever you
% want to typeset, there is probably a package out there for it.
% Chemistry (chemtex), screenplays, you name it.
% Check out CTAN to see: http://www.ctan.org/
%%
\usepackage{graphicx,latexsym}
% \usepackage{amsmath}
\usepackage{amssymb,amsthm}
\usepackage{longtable,booktabs,setspace}
\usepackage{chemarr} %% Useful for one reaction arrow, useless if you're not a chem major
\usepackage[hyphens]{url}
% Added by CII
\usepackage{hyperref}
\usepackage{lmodern}
\usepackage{float}
\usepackage{tikz}
\usetikzlibrary{shapes}
\usetikzlibrary{arrows}
\floatplacement{figure}{t}
% End of CII addition
\usepackage{rotating}

% Next line commented out by CII
%%% \usepackage{natbib}
% Comment out the natbib line above and uncomment the following two lines to use the new
% biblatex-chicago style, for Chicago A. Also make some changes at the end where the
% bibliography is included.
%\usepackage{biblatex-chicago}
%\bibliography{thesis}

\newlength{\cslhangindent}
\setlength{\cslhangindent}{1.5em}
\newlength{\csllabelwidth}
\setlength{\csllabelwidth}{3em}
\newlength{\cslentryspacingunit} % times entry-spacing
\setlength{\cslentryspacingunit}{\parskip}
% Version avec arguments (pour les nouvelles versions de pandoc)
\newenvironment{CSLReferences}[2] % #1 hanging-ident, #2 entry spacing
 {% don't indent paragraphs
  \setlength{\parindent}{0pt}
  % turn on hanging indent if param 1 is 1
  \ifodd #1
  \let\oldpar\par
  \def\par{\hangindent=\cslhangindent\oldpar}
  \fi
  % set entry spacing
  \setlength{\parskip}{#2\cslentryspacingunit}
  % Redéfinir \bibitem pour qu'il fonctionne hors d'un environnement liste
  \renewcommand{\bibitem}[2][]{\par}
 }%
 {}
% Version sans arguments (pour pandoc-citeproc)
\newenvironment{cslreferences}
 {\setlength{\parindent}{0pt}%
  \everypar{\setlength{\hangindent}{\cslhangindent}}\ignorespaces%
  \renewcommand{\bibitem}[2][]{\par}}
 {\par}
\usepackage{calc}
\newcommand{\CSLBlock}[1]{#1\hfill\break}
\newcommand{\CSLLeftMargin}[1]{\parbox[t]{\csllabelwidth}{#1}}
\newcommand{\CSLRightInline}[1]{\parbox[t]{\linewidth - \csllabelwidth}{#1}\break}
\newcommand{\CSLIndent}[1]{\hspace{\cslhangindent}#1}
\newcommand{\citeproctext}{}

% Added by CII (Thanks, Hadley!)
% Use ref for internal links
\renewcommand{\hyperref}[2][???]{\autoref{#1}}
\def\chapterautorefname{Chapter}
\def\sectionautorefname{Section}
\def\subsectionautorefname{Subsection}
% End of CII addition

% Added by CII
\usepackage{caption}
\captionsetup{width=5in}
% End of CII addition

% \usepackage{times} % other fonts are available like times, bookman, charter, palatino

% Syntax highlighting #22

% To pass between YAML and LaTeX the dollar signs are added by CII
\titre{Titre principal de votre document\\
Sous-titre (optionnel)}
\auteur{Prénom Nom}
% The month and year that you submit your FINAL draft TO THE LIBRARY (May or December)
\date{Mois Année}
\directeur{Prénom Nom du directeur}
\direction{directeur}
\codirection{}
\grade{M.A.}
\diplome{Maîtrise}
\type{Mémoire}
%If you have two advisors for some reason, you can use the following
% Uncommented out by CII
% End of CII addition

%%% Remember to use the correct department!
\departement{science politique}

% Added by CII
%%% Copied from knitr
%% maxwidth is the original width if it's less than linewidth
%% otherwise use linewidth (to make sure the graphics do not exceed the margin)
\makeatletter
\def\maxwidth{ %
  \ifdim\Gin@nat@width>\linewidth
    \linewidth
  \else
    \Gin@nat@width
  \fi
}
\makeatother

\renewcommand{\contentsname}{Table des matières}
% End of CII addition

\setlength{\parskip}{0pt}

% Added by CII

\providecommand{\tightlist}{%
  \setlength{\itemsep}{0pt}\setlength{\parskip}{0pt}}

\Remerciements{
Rédigez ici vos remerciements.
}

\Dedication{

}

\Preface{

}

\Resume{
Rédigez ici votre résumé (en français)
}

\Abstract{
Rédigez ici votre abstract (en anglais).
}

	\usepackage{setspace}
 \usepackage[labelfont=bf,textfont=md]{caption}
 \usepackage{float}
 \usepackage[nottoc]{tocbibind}
 \AtBeginDocument{\renewcommand{\contentsname}{Table des matières}}
 \AtBeginDocument{\renewcommand{\listfigurename}{Liste des figures}}
 \AtBeginDocument{\renewcommand{\listtablename}{Liste des tableaux}}
 \AtBeginDocument{\renewcommand{\chaptername}{Article}}
 \usepackage{pdflscape}
 \newcommand{\blandscape}{\begin{landscape}}
 \newcommand{\elandscape}{\end{landscape}}
% End of CII addition
%%
%% End Preamble
%%
%

\begin{document}

% Everything below added by CII
  \maketitle

\frontmatter % this stuff will be roman-numbered
\pagestyle{plain} % this removes page numbers from the frontmatter
\setcounter{page}{2}

  \begin{resume}
    Rédigez ici votre résumé (en français)
  \end{resume}

  \begin{abstract}
    Rédigez ici votre abstract (en anglais).
  \end{abstract}

% 
  \hypersetup{linkcolor=black}
   \setcounter{tocdepth}{2}
  \tableofcontents

  \listoftables

  \listoffigures

  \begin{remerciements}
    Rédigez ici vos remerciements.
  \end{remerciements}

  \begin{avantpropos}
    \hypertarget{duxe9claration-sur-lutilisation-de-lintelligence-artificielle-guxe9nuxe9rative}{%
    \section*{Déclaration sur l'utilisation de l'intelligence artificielle générative}\label{duxe9claration-sur-lutilisation-de-lintelligence-artificielle-guxe9nuxe9rative}}
    \addcontentsline{toc}{section}{Déclaration sur l'utilisation de l'intelligence artificielle générative}

    Avec l'accord de {[}ma directrice/mon directeur{]} de recherche, des outils d'intelligence artificielle générative {[}ont été utilisés/n'ont pas été utilisés{]} dans la préparation et la rédaction de ce {[}mémoire/cette thèse/cet essai{]}.

    Pour de l'information sur l'utilisation responsable de l'IA générative dans les mémoires et les thèses, consultez la page de la FESP : \url{https://www.fesp.ulaval.ca/memoires-et-theses/utilisation-responsable-de-lintelligence-artificielle-generative}
  \end{avantpropos}

% 
\mainmatter % here the regular arabic numbering starts
\pagestyle{fancyplain} % turns page numbering back on

\onehalfspacing

\hypertarget{introduction}{%
\chapter*{Introduction}\label{introduction}}
\addcontentsline{toc}{chapter}{Introduction}

Rédiger ici l'introduction de votre mémoire.

\hypertarget{votre-titre}{%
\chapter{Votre titre}\label{votre-titre}}

\hypertarget{abstract}{%
\section{Abstract}\label{abstract}}

\begin{quote}
Rédigez ici votre abstract (en anglais).
\end{quote}

\hypertarget{ruxe9sumuxe9}{%
\section{Résumé}\label{ruxe9sumuxe9}}

\begin{quote}
Rédigez ici votre résumé (en français)
\end{quote}

\begin{quote}
\textbf{Keywords}
Écrire ici vos mots-clés
\end{quote}

\hypertarget{texte}{%
\section{Texte}\label{texte}}

Rédigez ici votre étude.

Vous pouvez citer vos références directement dans le texte (Lewis-Beck, Norpoth, Jacoby, \& Weisberg, 2008). Les références peuvent inclure le numéro de la page (Alford, Funk, \& Hibbing, 2005, p. 13) ainsi que plusieurs références combinées (Converse, 2006; Iyengar \& Kinder, 2010; Page \& Shapiro, 2010).

\hypertarget{references}{%
\section{References}\label{references}}

\hypertarget{appendix}{%
\section{Appendix}\label{appendix}}

\hypertarget{votre-titre-1}{%
\chapter{Votre titre}\label{votre-titre-1}}

\hypertarget{abstract-1}{%
\section{Abstract}\label{abstract-1}}

\begin{quote}
Rédigez ici votre abstract (en anglais).
\end{quote}

\hypertarget{ruxe9sumuxe9-1}{%
\section{Résumé}\label{ruxe9sumuxe9-1}}

\begin{quote}
Rédigez ici votre résumé (en français)
\end{quote}

\begin{quote}
\textbf{Keywords}
Écrire ici vos mots-clés
\end{quote}

\hypertarget{texte-1}{%
\section{Texte}\label{texte-1}}

Rédigez ici votre étude.

\hypertarget{references-1}{%
\section{References}\label{references-1}}

\hypertarget{appendix-1}{%
\section{Appendix}\label{appendix-1}}

\hypertarget{conclusion}{%
\chapter*{Conclusion}\label{conclusion}}
\addcontentsline{toc}{chapter}{Conclusion}

Rédigez ici votre conclusion.

\backmatter

\hypertarget{bibliographie}{%
\chapter*{Bibliographie}\label{bibliographie}}
\addcontentsline{toc}{chapter}{Bibliographie}

\markboth{References}{References}

\hypertarget{refs}{}
\begin{cslreferences}
\leavevmode\hypertarget{ref-alford2005political}{}%
Alford, J. R., Funk, C. L., \& Hibbing, J. R. (2005). Are political orientations genetically transmitted? \emph{American Political Science Review}, \emph{99}(2), 153--167.

\leavevmode\hypertarget{ref-converse2006nature}{}%
Converse, P. E. (2006). The nature of belief systems in mass publics (1964). \emph{Critical Review}, \emph{18}(1-3), 1--74.

\leavevmode\hypertarget{ref-iyengar2010news}{}%
Iyengar, S., \& Kinder, D. R. (2010). \emph{News that matters: Television and american opinion}. University of Chicago Press.

\leavevmode\hypertarget{ref-lewis2008american}{}%
Lewis-Beck, M. S., Norpoth, H., Jacoby, W. G., \& Weisberg, H. F. (2008). \emph{The american voter revisited}. University of Michigan Press.

\leavevmode\hypertarget{ref-page2010rational}{}%
Page, B. I., \& Shapiro, R. Y. (2010). \emph{The rational public: Fifty years of trends in americans' policy preferences}. University of Chicago Press.
\end{cslreferences}


% Index?

\end{document}
